\answer{ex:07:1}{$V=\kappa p/\bigl(\kappa p+(1-\kappa)(1-p)\bigr)$: $0.059$, $0.20$, $0.50$; медіана дорівнює $0$, а точка~\eqref{eq:expectile} змінюється з $\kappa$ плавно.}
\answer{ex:07:2}{$p=1/3$, $x=0.75$, $\sigma_r=0.35$, $V(0.2)=0.083$.}
\answer{ex:07:3}{$V=\sqrt\kappa/(\sqrt\kappa+\sqrt{1-\kappa})$, від $0.25$ до $0.75$; розкид $\propto\sqrt\eta$; у крапель $V=0.25+0.5\kappa$, розподіли різняться на $0.05$ у крайніх нейронів, потрібно $\eta\lesssim0.01$.}
\answer{ex:07:4}{(а)~$J^*=2\sqrt{C\bar R}$. (б)~Переплата $\bigl(k^{1/2}+k^{-1/2}\bigr)/2-1$: $6\,\%$ при $k=2$ і $25\,\%$ при $k=4$ --- досить точності «в межах двох разів».}
\answer{ex:07:5}{Сплеск площею $R(e^{-1}-e^{-2})\approx0.23R$ --- як $2.3\,\text{с}$ рампи; якщо дофамін повідомляє $V$, сплеску немає, рівень просто сходинкою зростає в $e$ разів.}
\answer{ex:07:6}{Подія зсуває вагу $\propto A\tau_f/(\tau_e+\tau_f)$, рівень дає помилку $s\tau_f$: $9\,\text{с}\le\tau_f\le17\,\text{с}$.}
\answer{ex:07:7}{Вихідно $\Delta P=0.80$, $M=0.90$. (а)~$\Delta P=0.74$ ($-8\,\%$), $M=0.43$ ($-52\,\%$); (б)~$\Delta P=0.40$ ($-50\,\%$), $M=0.80$ ($-11\,\%$): подвійна дисоціація.}
\answer{ex:07:8}{$N_{\text{еф}}\to2+\rho^2N\approx10^6$ спроб: у $10^4$ разів менше, ніж з одним проводом, але витік в $1\,\%$ усе ще коштує мільйона спроб.}
